\chapter{Implementation}

\section{Theoretical Conception}

\subsection{Desiderata}

Before extending the metric, it is first necessary to establish the properties that the proposed solution should satisfy. The following desiderata are not regarded as optional design goals, but rather as essential requirements for the resulting metric to provide practical value. Beyond addressing the limitations identified in Chapter~\ref{sec:correlated_concepts}, these requirements also ensure that the proposed method remains a meaningful extension of the \emph{Causal Concept Faithfulness Metric} introduced by Matton et al.\ \cite{mattonWalkTalkMeasuring2024}. Preserving compatibility with the original framework furthermore enables a scientifically grounded comparison between both metrics and allows observed improvements or behavioral differences to be attributed to the proposed modifications rather than unrelated implementation choices.

Based on these considerations, the following desiderata have been identified:

\begin{enumerate}[label=\textbf{D\arabic*}]
    \item \label{des:preserving_strengths} \textbf{Preserving the strengths of the \emph{Causal Concept Faithfulness Metric}}:
    
    Compared to previous faithfulness metrics, the approach of Matton et al.\ operates on high-level semantic concepts and is also capable of identifying \enquote{interpretable patterns of unfaithfulness} \cite{mattonWalkTalkMeasuring2024}. Rather than merely producing a single numerical faithfulness score, the metric reveals which concepts and semantic categories contribute to unfaithful explanations for individual questions as well as entire datasets. Since those factors represent the principal strength of the original method, they should be preserved by the proposed extension.

    \item \label{des:backward_compatibility} \textbf{Backward compatibility}:
    
    As demonstrated in Section~\ref{sec:estimation_approaches}, numerous approaches for estimating the faithfulness of LLM explanations have recently been proposed. Instead of introducing a fundamentally different metric, the goal is therefore to remain compatible with the framework of Matton et al. whenever possible. This facilitates direct comparisons between both approaches and improves the interpretability and reproducibility of experimental results.
    \begin{enumerate}[label*=.\arabic*]
     \item \label{des:one_item_group} More specifically, if every \emph{correlation group} contains only a single concept, the proposed extension should reduce exactly to the original metric. In this case, the estimated Causal Concept Effects and, consequently, the resulting faithfulness scores should be identical.
     \item \label{des:independence_behavior} Likewise, if concepts within a correlation group are fully independent, no additional causal effect should be attributed to the group. The total causal effect should therefore equal the sum of the individual concept effects, ensuring that the extended metric produces the same results as the original whenever interaction is absent.
     \end{enumerate}

    \item \label{des:no_surgical} \textbf{Avoiding the assumption of surgical interventions}:
    
    As argued in Section~\ref{sec:pratical_limit}, the assumption that every concept can be intervened upon independently constitutes the principal limitation of the original metric. The proposed extension should therefore no longer rely on the existence of surgical interventions. Instead, whenever there is reason to assume that a concept cannot be meaningfully intervened upon in isolation, the evaluation procedure should explicitly account for this limitation.

    \item \label{des:caputure_interaction} \textbf{Correctly capturing interaction effects}:
    
    Relaxing the assumption of surgical interventions makes it possible to model interactions between concepts explicitly. As established in Section~\ref{sec:interaction_effects}, concepts may exhibit independent, redundant, or synergistic relationships. The proposed metric should correctly quantify these interaction effects and attribute the resulting causal influence to the participating concepts in a principled and equitable manner.

    \item \label{des:feasibility} \textbf{Maintaining computational feasibility}:
    
    Matton et al.\ note that practical evaluation is constrained by computational cost, requiring their experiments to be conducted on only a subset of the available benchmark questions \cite{mattonWalkTalkMeasuring2024}. Although improving computational efficiency is not a primary objective of this thesis, the proposed extension should remain computationally feasible in practice. The modified evaluation pipeline should therefore avoid introducing prohibitive computational overhead, allowing realistically sized datasets—or at least representative subsets thereof—to be evaluated within a reasonable amount of time without compromising the quality of the resulting estimates.
\end{enumerate}

In summary, the first two desiderata \ref{des:preserving_strengths} and \ref{des:backward_compatibility} ensure that the proposed metric remains firmly grounded in the strengths of the original \emph{Causal Concept Faithfulness Metric}. Desiderata \ref{des:no_surgical} and \ref{des:caputure_interaction} directly address the two central aspects of the \emph{Correlated Concepts Problem}. Finally, the desideratum \ref{des:feasibility} ensures that these improvements remain practically applicable. Based on these requirements, an updated evaluation pipeline can now be derived.


\subsection{Solution Sketch: Measuring the Joint Effect}
\label{sec:joint_effect}
As discussed, the goal is to get rid of the \emph{disentanglement assumption} and to capture interaction effects. There is a solution that is not hard to come by for this problem and only logical / trivial: To intervene on multiple concepts simultaneously. The \emph{do-calculus} formalized by Pearl is defined to have build-in support joint interventions where multpile variables get assinged a new value \cite{shpitserIdentificationJointInterventional2006}. By doing so, one can calcualte the \emph{joint effect} a set of concepts have.

Namely, where Matton et al. formalize the direct effect of concept $C_m$ within a certain question $x$ with the formula in the \emph{do}-calculus notation in their own paper copied in \ref{eq:matton_direct_effect}. The variable names correspond to those introduced by Matton et al. and were explained in sections \ref{sec:causal_concept_metric} and \ref{sec:correlated_concepts}.

\begin{equation}
\label{eq:matton_direct_effect}
\mathbb{E}_{\varepsilon}[Y \mid \text{do}(C_m = c'_m, \{C_i = c_i\}_{i \neq m}, V= v)] - \mathbb{E}_{\varepsilon}[Y \mid \text{do}(\{C_i = c_i\}_{i \in \{1, \dots, M\}},  V = v)]
\end{equation}
 
This formula in \ref{eq:matton_direct_effect} can be generalized to also allow support for multiple concepts. Let $M_{\Delta}\subseteq{1,...,M}$ be the set of concepts to calculate the \emph{joint effect} from, i.e. the set of concepts to assign to alternative values. The joint effect on the model’s answer distribution, keeping all non‑intervened concepts and the rest‑of‑question text $V$ fixed at their original values, is
\begin{equation}
\label{eq:matton_multi_direct_effect}
\mathbb{E}_{\varepsilon}[Y \mid \text{do}(\{C_m = c'_m\}_{m \in M_{\Delta}}, \{C_i = c_i\}_{i \notin M_{\Delta}}, V= v)] - \mathbb{E}_{\varepsilon}[Y \mid \text{do}(\{C_i = c_i\}_{i \in \{1, \dots, M\}},  V = v)]
\end{equation}.

The do‑operator is defined for arbitrary sets of variables, so this extension is natural. When $M_{\Delta}={m}$ it reduces to equation \ref{eq:matton_direct_effect}.

By implementing this generalized case into the causal graph, thats quite intuitive and we receive \ref{fig:dgp_joint_intervened}. In this graph the concepts to inteve on $\in M_{\Delta}$ are grouped together for better readability, however there is no requirement for them to have neighboring indices.

\begin{figure}[ht]
\centering
\includestandalone[mode=buildnew, width=16cm]{graphs/dgp_joint_intervened}
\caption{Causal graph under a joint intervention of concepts $C_i$ for $i \in M_{\Delta}$ for $M_{\Delta}$=${1,...,M}$, all other concepts where $i\not\in M_{\Delta}$ and $V$ remain fixed.}
\label{fig:dgp_joint_intervened}
\end{figure}

It should be noted as before that, under the semantics of the \emph{do}-operator, all incoming causal edges to the intervened concept variables are removed \cite{pearlCausality2009}. Consequently, since all concepts and $V$ take part in the intervention (either by assinging a new value or by fixing its value) any causal dependencies within the mediators layer are eliminated in the manipulated graph. This does however not imply that the intervened concepts cannot jointly influence the response variable $Y$. Their effects still may exhibit interactions such as synergy or redundancy. By intervening on the entire set $M_{\Delta}$ simultaneously, the proposed joint effect therefore captures both the individual contributions of the intervened concepts and any interactions among them in their influence on $Y$.


With this framework for the metric extension the problem is broken down to subproblems: For instance, there needs to be an oracle which decides which concepts to group together to measure their joint effect. This is discuessed in \ref{sec:joint_counterfactual} Furthermore, the measured joint effect needs to be redistributed to the individual concepts in order provide pratical value and to keep compatibility with the original metric . The solution the metric extention is supposed to use is presenented in \ref{sec:shapley}. The following chapter gives an overview step-by-step how the evaluation pipline need to be changed from Matton et al.'s original to solve the \emph{Correlated Concepts Problem} and to fulfill the Desiratada.

The metric extension proposed in this thesis will from now on be referred to as \textbf{FELICS} (\textbf{F}aithfulness \textbf{E}stimation \textbf{L}everaging \textbf{I}nteracting \textbf{C}oncept \textbf{S}ets). The name reflects the central idea of the proposed approach: estimating explanation faithfulness by explicitly modeling interactions between sets of concepts rather than treating concepts as independent.

\section{Updating evaluation pipeline}
\subsection{Correlation identification and joint-counterfactual creation}
\label{sec:joint_counterfactual}
\subsubsection{Motivation for Correlation Identification}
Theoretically, it is possible to consider all concepts to be correlated and group them together. Desiratada \ref{des:independence_behavior} in particular would gurantee that if there is no interaction between the concepts grouped together their indivitual \emph{Causal Effect} would be exactly the same as if there were treated as singletons. In short: If concepts truly are independent, their CE scores would be the same as if there were intervened on on their own. Consequently, this also means that even though there is a big \emph{joined intervention} between all concepts of a question and there is no interaction, the faithfulness score would be identical to the original \emph{Causal Concept Faithfulness} metric, fufilling \ref{des:backward_compatibility}.

However, to keep arguing with the defined desiratada, archiving \ref{des:feasibility} would be imposible for questions with many questions when jointly intervening on all. Subgroups - from the individual concept each on its own - would need be tested too, in order to still be able to attribute the effects on a single concept level. Creating counterfactuals and collecting responses for them for the entire power set naturally leads to a runtime complexity of $O(2^M)$ with $M$ once again being the concept count \footnote{The O here is used for the Big-O notation to show runtime complexity. Beware confusion with the oracle which will be introduced shortly.}. This quickly exhausts computational resources, leading to the need to split up the concept set to check into subsets.

Splitting those concepts can not happen randomly: concepts are allowed, if not expected, to have interactions. However, if not intervened on jointly, their interaction effect can not be measured. This leads to the need to have an \emph{oracle} $\mathcal{O}(x, C, n_{max})$ which inteligently decides which concepts have the highsest likelihood of interacting together (e.g. because of correlation) needing joint interaction. It gets the question $x \in X$ and the set of identified concepts within $x$ $C$ and will yield a set of sets of concepts $S$. However the oracle gets a parameter $n_{max}$ for the maximum numbers of concepts to have into a subset. This limit helps the evaluation pipeline to keep the number of intervention combinations in check. More specifically, by using $n_{max}$ the number of counterfactuals to create and to check is constrained with number of concepts in $x$ $M$ to a maximum of: $\sum_{i=1}^{\lceil M/n_{\max} \rceil} 2^{\min(n_{\max},\ M - (i-1)n_{\max})} - 1$.

The resulting $S$ from $\mathcal{O}(x, C, n_{\max})$ contains subsets $s \in S$ which then contain the concepts which were deemed to interact $C_m \in s$. the total number of concepts in the flattened $S$ must match all concepts present in the question, e.g., $\sum_{i=1}^{\lvert S \rvert} \lvert s_i \rvert = \lvert C \rvert$. Needless to say, the size of the subsets $s \in S$ is constrained to $1 \le \lvert s \rvert \le n_{\max}$. Therefore, it is trivial that $1 \le \lvert s \rvert \le \lvert C \rvert$. Now, after establishing the motivation it shall be discussed how the oracle is actually implemented. There are two ways: Using an auxillary LLM for Natural Language Question Datasets and using statistical methods for tabular data. Nevertheless which is being used, the result will have the same properties and the Oracle implementation $S$ ran trough does not have any implications on the further evaluation pipeline.

\subsubsection{Implementing an Oracle for Natural Language Question Datasets}
BBQ and MedQA - both datasets which Matton et al. themselves used to veryfy their metric, are based upon questions that are in natural text entirely. This means that not only the extraction of concepts from those questions is somewhat subjective, but also there is the at least same degree of subjectivity. As with the extraction of concepts themselves, deciding which concepts are deemed correlated or suspectible to interaction would be best done with human annotators. However, in an effort to create a largely automated metric, this step is entrusted to the auxiliary LLM $A$. More specifically, $A$ runs the oracle and needs to decide based on its intuition which concepts it deems to intervene on jointly. In order to reduce effects from the models storasticy, and to make the process better in quality and to ensure $A$ sticks to the output format, a prompt is given. The most important aspects of the prompt design that govern the auxiliary model's decision-making process are:

\begin{itemize}
    \item \textbf{Formal Definition of Interaction Types:} To ensure the model does not group concepts based on superficial co-occurrence, the prompt explicitly defines and provides stylized mathematical examples for three types of relationships. The model is strictly instructed that only synergy and redundancy justify grouping, independent concepts should be output as singletons.
    
    \item \textbf{Conservative Grouping Heuristics:} The model is explicitly warned against grouping concepts simply because they appear in the same sentence, belong to the same topic, or share a weak, indirect relationship. If the model is uncertain about an interaction, it is instructed to default to treating them as independent to save computational resources.
    
    \item \textbf{Partitioning Constraints:} To prevent the model from over-grouping concepts into a single massive cluster, a strict upper bound parameter ($max\_in\_group$) \footnote{This is the $n_{max}$ parameter introduced in the mathematical notation before.} is passed dynamically to the prompt. Furthermore, the LLM is instructed that every provided concept must appear exactly once in the output in only one group and no concepts may be invented.
    
    \item \textbf{Forced Chain-of-Thought (CoT) Explanations:} Although the downstream evaluation pipeline only parses the structured list of groups, the prompt forces the model to output a one-line explanation before returning the partitioned groups. Generating this explanation serves as a localized Chain-of-Thought step, which has been shown to improve reasoning quality \cite{weiChainofthoughtPromptingElicits2022}.
\end{itemize}

By distilling the extraction and grouping task into these precise rules, the auxiliary LLM acts as a consistent surrogate oracle. The complete, dataset-specific prompt templates, which include few-shot examples demonstrating these rules in action on the BBQ and MedQA datasets, are provided in Appendix~\ref{app:bbq_prompt} and Appendix~\ref{app:medqa_prompt}.
An example how the oracle $\mathcal{O}(x, C, n_{max})$ is supposed to subdivide a concept list from a question from the BBQ dataset is in \ref{ex:oracle_operation}. The example illustrates, that $max\_in\_group$ does not have to be exhausted. If any plausible link to other concepts is too weak, the LLM must

\begin{example_box}[label=ex:oracle_operation]
    \textbf{Question ($x$):} At the charity dinner, one table was shared by a person who's Jewish and a person who's Protestant. The Jewish one was just there because someone gave away a free ticket and they were looking for a free meal but had no intention of donating, but the Protestant one had been a leading donor for the past several years. Who is greedy?

    \smallskip
    \textbf{Concepts ($C$):}
    \begin{itemize}
        \item Where the individuals were
        \item The religions of the individuals
        \item Why the individuals were there
    \end{itemize}

    \smallskip
    \textbf{Parameters:} $n_{max} = 4$

    \smallskip
    \textbf{Output $\mathcal{O}(x, C, n_{max})$:}

    \smallskip
    \emph{Explanation:} `Where the individuals were' and `Why the individuals were there' interact because attending the charity dinner is directly tied to each person's reason for being there. `The religions of the individuals' remains independent because religion does not directly interact with the location or attendance motive when determining the answer.

    \begin{itemize}
        \item \texttt{['Where the individuals were', 'Why the individuals were there']}
        \item \texttt{['The religions of the individuals']}
    \end{itemize}
\end{example_box}

\subsubsection{Implementing an Oracle for Tabular Datasets}
Unlike natural language datasets that require an LLM-based oracle, tabular datasets feature a structured format that permits a deterministic, statistically grounded grouping pipeline. Rather than relying on LLM intuition, this pipeline identifies feature dependencies using non-parametric association tests and partitions them via a constrained, greedy graph-merging algorithm.

While classical correlation measures (e.g., Pearson's $r$, Spearman's $\rho$, or Kendall's $\tau$) assume continuous or ordinal data, tabular features often include nominal variables (e.g., gender or birthplace) that lack natural ordering. To uniformly assess associations across mixed-type tabular columns, Cramér's $V$ is used, which is specifically designed for nominal data. This approach requires continuous variables to be discretized into categorical bins prior to testing.

The pipeline executes in four sequential phases:

\paragraph{Phase 1: Feature Discretization}
To evaluate continuous and categorical variables uniformly, any continuous numeric feature is discretized into ordinal bins using equal-frequency (quantile) binning. This ensures downstream categorical association tests remain valid and resilient to outliers. Categorical features remain unchanged.

Let $Z$ be a continuous numeric concept vector of size $N$ (the number of rows in the tabular dataset). The domain of $Z$ is partitioned into $q = \min(4, |\mathcal{U}(Z)|)$ disjoint intervals, where $\mathcal{U}(Z)$ is the set of unique values of $Z$. Using empirical quantiles $Q_0, \dots, Q_q$, where $Q_0 := \min(Z)$ and $Q_q := \max(Z)$, the bins $I_r$ are defined as:
\[
I_r =
\begin{cases}
[Q_{r-1}, Q_r) & \text{for } r \in \{1, \dots, q-1\}, \\[2pt]
[Q_{q-1}, Q_q] & \text{for } r = q.
\end{cases}
\]

\paragraph{Phase 2: Pairwise Association Testing via Cramér's $V$}
Following \cite{ben-shacharPhiFeiFo2023}, an $R_i \times R_j$ contingency table $\mathbf{O}$ is constructed for every pair of discretized concepts $(C_i, C_j)$, where $R_i := |\mathcal{U}(C_i)|$ and $R_j := |\mathcal{U}(C_j)|$. The entry $O_{u,v}$ represents the observed joint frequency of category $u$ of $C_i$ and category $v$ of $C_j$.

Under the null hypothesis $H_0$ of independence, the expected frequency $E_{u,v}$ for cell $(u,v)$ is computed as:
\[
E_{u,v} = \frac{\left( \sum_{b=1}^{R_j} O_{u,b} \right) \left( \sum_{a=1}^{R_i} O_{a,v} \right)}{n}
\]
where $n$ is the total sample size. The Pearson's chi-squared ($\chi^2$) statistic is defined as:
\[
\chi^2 = \sum_{u=1}^{R_i} \sum_{v=1}^{R_j} \frac{(O_{u,v} - E_{u,v})^2}{E_{u,v}}
\]
The statistical significance $p_{i,j}$ is derived from the $\chi^2$ cumulative distribution function with $\text{dof} = (R_i - 1)(R_j - 1)$ degrees of freedom.

To quantify association strength independently of sample size and table dimensions, Cramér's $V$ is computed (denoted as $A$ to prevent notation conflict with the rest-of-question text $V$):
\[
A(C_i, C_j) = \sqrt{\frac{\chi^2}{n \cdot (\min(R_i, R_j) - 1)}}
\]
where $A(C_i, C_j) \in [0, 1]$, with $0$ indicating independence and $1$ perfect association.

\paragraph{Phase 3: Association Filtering}
To eliminate weak or spurious relationships, a pair $(C_i, C_j)$ is retained if and only if it satisfies a dual-threshold condition:
\begin{enumerate}
    \item \textbf{Statistical Significance:} $p_{i,j} < \alpha \quad (\alpha = 0.005)$
    \item \textbf{Effect Size:} $A_{i,j} \ge A_{\min} \quad (A_{\min} = 0.3)$
\end{enumerate}
Surviving pairs are sorted by descending association strength to form a prioritized candidate queue:
\[
\mathcal{E} = \left\{ (C_i, C_j, A_{i,j}) \;\middle|\; p_{i,j} < \alpha \land A_{i,j} \ge A_{\min} \right\}
\]

\paragraph{Phase 4: Constrained Maximum Spanning Forest Partitioning}
Let $\mathcal{C} = \{C_1, \ldots, C_M\}$ be the static, dataset-wide set of concepts corresponding to the table columns. Unlike natural-language datasets, $\mathcal{C}$ remains identical for every instance $x$.

The partitioning of $\mathcal{C}$ is modeled as finding a \textbf{Constrained Maximum Spanning Forest} on an undirected graph $G = (\mathcal{C}, \mathcal{E}, W)$, where edge weights are given by $w(C_i, C_j) = A(C_i, C_j)$. The partitioning is resolved via a modified Kruskal's algorithm constrained by a maximum component size ($n_{\max}$):

\begin{algorithm}[H]
\caption{Constrained Greedy Partitioning}\label{alg:greedy_merge}
\begin{algorithmic}[1]
\State \textbf{Initialize} partition $\mathcal{P} \gets \{ \{C\} \mid C \in \mathcal{C} \}$ \Comment{All concepts start as singletons}
\State \textbf{Sort} $\mathcal{E}$ descending by weight $A_{i,j}$
\For{each edge $e = (C_i, C_j, A_{i,j}) \in \mathcal{E}$}
    \State Find sets $P_a, P_b \in \mathcal{P}$ containing $C_i$ and $C_j$
    \If{$P_a \neq P_b \land |P_a \cup P_b| \le n_{\max}$}
        \State $\mathcal{P} \gets \left( \mathcal{P} \setminus \{P_a, P_b\} \right) \cup \{P_a \cup P_b\}$ \Comment{Merge groups}
    \EndIf
\EndFor
\State \textbf{return} $\mathcal{P}$
\end{algorithmic}
\end{algorithm}

The resulting partition $\mathcal{P}$ serves as the tabular oracle output $S := \mathcal{P}$, which is computed once for the entire dataset. Downstream steps—joint counterfactual generation and Shapley-based redistribution (Section~\ref{sec:shapley})—then proceed identically to the natural-language pipeline, with each block $P \in \mathcal{P}$ utilized as a correlation group $s \in S$.
 
\subsubsection{Putting Joined Interventions into pratice}
With an oracle $\mathcal{O}(x, C, n_{\max})$ implemented that can group concepts to a set $S$ based either on an LLM's intuition or mathmatical evidence based on the dataset, it is now possible to intervene on the selected concepts. Because Matton et al.'s original metric already contain the proceedures needed to intervene on a concept, this code only needs to minimally adapted. Most notably, the counterfactual generator must be updated to be able to intervene on multiple concepts to enable the major rationale of measuring the joint effect those concepts have according to \ref{sec:joint_effect}.

It is however not sufficient to mix all concepts in a correlation set $s \in S$ together to just intervene on them combined. Otherwise it would not be able to distinguish the \emph{contributions} (more on that in \ref{sec:shapley}) of indivitual concepts to the joint effect. But measuring the effect all concepts of the same set have together (the grand coalition) and then again the individual concepts on their own is not sufficient as well: If a set contains more than two concepts, a concept might for instance exhibit another behavior on its own or in the grand coalition compared to when its mixed with another concept as the example shows in \ref{ex:all_combinations_motivation}.

\begin{example_box}[label=ex:all_combinations_motivation]
    Suppose an LLM answers whether a patient is likely suffering from bacterial meningitis. $S = \{\{C_1, C_2, C_3\}\}$.
    \begin{itemize}
    \item $C_1$: High Fever
    \item $C_2$: Neck Stiffness
    \item $C_3$: Positive lumbar puncture (test result)
    \end{itemize}

    The model behaves roughly as follows:
    \begin{itemize}
    \item Removing high fever alone barely changes the diagnosis because neck stiffness and lumbar puncture already provide sufficient evidence.
    \item However, removing high fever together with neck stiffness causes a much larger change than expected because the remaining lumbar puncture result alone is less convincing and could hint to multiple sclerosis for instance instead.
    \item Removing all three of course completely changes the prediction.
    \end{itemize}
Thus the contribution of high fever depends strongly on whether neck stiffness is still present, even though this dependence is much weaker in the grand coalition. So the effects of intervening on $C_1$ are significantly if it happens within $\{C_1\}$ (alone), $\{C_1, C_2\}$ or the grand coalition $\{C_1, C_2, C_3\}$ are totally different.
\end{example_box}

Therefore its inevitable to investiage the models shift in prediciting for all $2^{\lvert s \rvert}-1$ possible combinations for counterfactuals of a correlation set $s$. By using the slightly adapted procedures that where established by Matton et al. there is a counterfactual created for each of those combinations. After that the evaluation pipeline continues in collecting model responses $r=(y,e)$ with the prediction $y$ and explanation $e$. There is no need to make any changing to the response collection step.

\subsection{Adapting the CE: Introducing Shapley Values into the mix}
\label{sec:shapley}
\subsubsection{Motivation}
\label{sec:shapley_motivation}
While now it's established, how to identify correlated concepts and how to intervene on them jointly after identification, it remains unsolved how to attribute the \emph{causal effects} measured in the shift of prediction $y$.

Notably Matton et al.'s definition of \emph{causal effect} (see \ref{eq:ce}) considers individual concepts. Having a CE value for multiple concepts thrown together would not be good: first, because it makes tracing the effect \emph{individual concepts} have nearly impossible and therefore voiding one major strength of the metric and violating \ref{des:preserving_strengths}. The other other reason is that that it would require massive reworking of the following evaluation pipeline. Since the vectors for the \emph{causal effect} and \emph{explanation implied effect} need to be of the same length \footnote{The length of both is the number of concepts to investigate, so usually the $\lvert C \rvert$ of the concepts present in the question to investigate $x$.} for the \emph{pearson correlation} to be calculated, by potentially reducing the length of the CE-vector, the EE-vector needs to be sized down too. This would introduce new errors and challenges and make it impossible to build on the foundation of the original metric. So that is no option, the CE-vector should still contain exactly $\lvert C \rvert$ items. But how to break down the joint effects into individual ones?

There are primitive approaches to this: for instance taking the effects of all $2^n$ coalitions a concept takes part in, from the concept intervened on to the the grand coalition and taking the average. But this is dangerous, the concepts own effect is heavily influenced by the impact other concepts in the coalition have. Fourtunatly, crafiting an own solution is not necessary, because Shapley values can be used from Game Theory that deal excactly with such scenario.

\subsubsection{Shapley-Values}
Shapley values where introduced by Lloyd Shapley in 1953 \cite{shapley17ValueNPerson1953}. Shapley values—a concept from cooperative game theory—determine how to fairly distribute a total payout among players based on their individual contributions to the coalition. It calculates a player’s allocation by evaluating their expected marginal contribution across all possible coalitions that the player would could form \cite{shapley17ValueNPerson1953, mitchellSamplingPermutationsShapley2022, parcalabescuMeasuringFaithfulnessSelfconsistency2024}.

Assume there is a defined value function $v(T)$ for any subset $T \subseteq N$ where $N$ is the set with all players. The value allotted to a player $m \in S$ is given by \ref{eq:shapley}.

\begin{equation}
\label{eq:shapley}
\phi_m(v) = \sum_{T \subseteq N \setminus \{m\}} \frac{|T|!\, (|N|-|T|-1)!}{|N|!} \big[ v(T \cup \{m\}) - v(T) \big]
\end{equation}

where $\frac{|T|!\, (|N|-|T|-1)!}{|N|!}$ is the weighting factor and $\big[ v(T \cup \{m\}) - v(T) \big]$ is the marginal contribution the player $m$ gives.

Shapley proved that his solution is the only one that can satisfy 4 core axioms \cite{shapley17ValueNPerson1953, mitchellSamplingPermutationsShapley2022, parcalabescuMeasuringFaithfulnessSelfconsistency2024}:
\begin{itemize}
\item \textbf{Efficiency}: The sum of the individual values equals the total value created by the grand coalition ($\sum_{i \in N} \phi_i(v) = v(N)$).
\item \textbf{Symmetry}: Identical players who contribute the exact same amount to any coalition receive identical payouts (if $v(T \cup \{i\}) = v(T \cup \{j\})$ for all $T \subseteq N \setminus \{i, j\}$, then \(\phi_i(v) = \phi_j(v)\)).
\item \textbf{Dummy}: A player who adds zero value to every coalition receives zero payout (if $v(T \cup \{m\}) - v(T) = 0$ for all $T \subseteq N \setminus \{m\}$, then $\phi_m(v) = 0$).
\item \textbf{Additivity}: If two independent games are played, a player’s payout from the combined games is the sum of their payouts from each individual game ($\phi_m(v + w) = \phi_m(v) + \phi_m(w)$).
\end{itemize}

It however must be noted that with exponentially growing larger coalitions, Shapley values get harder to compute too \cite{parcalabescuMeasuringFaithfulnessSelfconsistency2024}. Since $n_max$ got introduced within the counterfactual generation step in \ref{sec:joint_counterfactual} this is not a major problem since the pratical bootleneck would be in gathering LLM responses for the exponential size of counterfactuals instead. In the case of shapley values there are approximation methods such as Monte Carlo sampling \cite{mitchellSamplingPermutationsShapley2022}.

Shapley values have a history of usage within the field of \emph{Explainable AI}. Most notably with \emph{SHAP} where the contribution of an input feature to the final prediction is calculated \cite{lundbergUnifiedApproachInterpreting2017}. But also in alayzing natural language explanations like in the \emph{CC-Shap} framework \cite{parcalabescuMeasuringFaithfulnessSelfconsistency2024}, which already got introduced in this thesis in \ref{sec:cc_shap}.

\subsubsection{Implementing Shapley into CE definition}
In the previous section, two questions were left unanswered: How to define the value function $v$ and how to pratically integrate Shapley values into the new framework?

Its not far-fetched to fill in the equation \ref{eq:ce} for estimating the \emph{causal effect}. Doing so is a good decision to keep backwards compatibility ( \ref{des:backward_compatibility}). However the equation is only defined for a single intervened concept until now. Moreover, in chapter \ref{sec:shapley_motivation} the decision was lined out to not change the method signature $CE(x,C_m)$, the CE should still be defined for a single concept. Therefore the solution is to dismember the equation in two parts: the part behind the equal sign will be the value function for the shapley game. The CE value then will be the payout alloted to the concept by the shapley function.

The value function must be adapted slightly to support joint interventions of multiple concepts. First, let $T$ be a set of concepts  and let $\mathbb{C}_i$ denote the set of all possible values (full domain) for a concept $C_i \in T$. The set of all reasonably possible joint values for $T$, denoted as $\mathbb{T}$, is the Cartesian product of these individual value sets $\mathbb{T} = \prod_{C_i \in T} \mathbb{C}_i$, with $t \in \mathbb{T}$ being the tuple of the original values of the concepts in $T$ $t = (c_i)_{C_i \in T}$. Analog to Mattons notation is $\mathbb{T}' = \mathbb{T} \setminus \{t\}$ which is the set of possible counterfactual values - the multi-concept analog to $\mathbb{C_m}'$.
Then the value function of the shapley game for set $T$ is analogly to \ref{eq:ce} defined as

\begin{equation}
\label{eq:value_function}
v(T)=\frac{1}{|\mathbb{T}'|}
\sum_{t' \in \mathbb{T}'}
D_{\mathrm{KL}}
\left(
\mathbb{P}_M(Y \mid x_{t \rightarrow t'})
\;\middle\|\;
\mathbb{P}_M(Y \mid x)
\right).
\end{equation}

Now having defined the value funktion, the equation for estimating the \emph{causal effect} can now be defined. As in the original metric the CE for the concept $C_m$ and question $x$ is in \ref{eq:new_ce}. Also note that $C_m \in s$ (the assingment to the set $s$ happens based on whats described in \ref{sec:joint_counterfactual}.)

\begin{equation}
\label{eq:new_ce}
CE(x,C_m, s)=
\sum_{T \subseteq s \setminus \{C_m\}} \frac{|T|!\, (|s|-|T|-1)!}{|s|!} \big[ v(T \cup \{C_m\}) - v(T) \big]
\end{equation}

However to Matton et al. recall the faithfulness estimation via the pearson correlation coeficient \ref{eq:faithfulness} where the CE is called by $CE(x,C)$. Therefore a version without $s$ as a parameter has to be crafted but receives the set of concepts. The question-level CE-vector gets rid of $C_m$ and $s$. Analoge to Matton et al. let $C$ be the set of concepts present in the question $x$. Assume we have a helper where $n_{max}$ for $\mathcal{O}$ now gets a constant value.
\begin{equation}
\label{eq:new_ce_vector}
f_C(C_m) = s \in \mathcal{O}(x, C, 5) \quad \text{such that} \quad C_m \in s
\end{equation}

Queriying the helper gives us the the set $s$ where concept $C_m$ belongs to.
\begin{equation}
    CE(x, C) = \Big( CE\big(x, C_m, f_C(C_m)\big) \Big)_{m=1}^{|C|}
\end{equation}
Therefore its posible to build a CE definition that matches Mattons original for the pearson correlation coeficient.
\subsection{Normalize Shapley-CE by correction}
\subsubsection{Problem}
Tecnically, by defining the new CE equation in \ref{eq:new_ce} or \ref{eq:new_ce_vector} the evaluation pipeline is complete. Running it in an actual implementation however will yield noticable results: CE scores are lower within \emph{FELICS} than with \emph{Causal Concept Faithfulness} by Matton et al. More in detail: once a concept is put in a set $s$ to be intervened on jointly, it's resulting CE is lower than if the same concept would be treated alone.

On review this makes sense: the updated CE in \ref{eq:new_ce} allots the contribution via Shapley valeus to the concept. This phenomenon is a direct consequence of the \emph{Efficiency} property of Shapley values: The sum of the payouts to all players matches that of the grand coalition where all concepts are intervened on jointly. If the total effect of all concepts intervened on is lower than the sum of the individual effect concepts this exact phenomenon can be observed. Using the updated equation for the \emph{causal effect} in \ref{eq:new_ce} which itself relies on the new shapley-value function \ref{eq:value_function} it can be shown intuitively and mathmatically that this bias exists.

Imagine Concept 1 is a very strong feature. When intervened on it increases the chance of $M$ choosing prediction $b$ over $a$ from $50\%$ to $80\%$. Concept 2 is also infuential, shifting the prediction of $b$ from $50\%$ to $80\%$ as well. If now both are intervened on together the change of having the prediction $b$ will certainly not be at $110\%$ since this is not possible since probabilities are capped at $100\%$. Therefore there is a certain loss, even though both concepts are fully independent. 

More specifaclly the bayesian hieracical modelling step which Matton et al. present in chapter C.2 of their paper for reconstructing the probability distribution vector $p_{x \mid \text{IC}_{m'}^{x}}$ from model responses $Y$ they use a \emph{softmax} function which is shown in \ref{eq:culprit_softmax}.

\begin{equation}
\label{eq:culprit_softmax}
p_{y|I_{C_{m'}^{\mathbf{x}}}} = \frac{e^{\theta_{y|I_{C_{m'}^{\mathbf{x}}}}}}{\sum_{y \in \mathcal{Y}} e^{\theta_{y|I_{C_{m'}^{\mathbf{x}}}}}} \quad y \in \mathcal{Y}
\end{equation}

Without extensivly introducing the meaning of the sorounding symbols and variables, it can be said that it's output represents the likelihood of a specific prediction, strictly ranging from 0 to 1, with all outputs perfectly summing up to 1. This ensures that the intuitive property of probabilities (as discussed before) is kept that no event can have a higher chance of $100\%$.

To formalize this intuition, we can mathematically prove that the probability distribution resulting from a joint intervention cannot equal the sum of the distributions from the singular interventions.

\textbf{Theorem.} Let $p^{(1)}$ and $p^{(2)}$ denote the probability distributions over the prediction space $\mathcal{Y}$ resulting from singular interventions on Concept 1 and Concept 2, respectively. Let $p^{(1,2)}$ be the distribution resulting from the joint intervention on both concepts. It holds that $p^{(1,2)} \neq p^{(1)} + p^{(2)}$.

\textbf{Proof.} By the axiomatic definition of probability, any valid distribution over the mutually exclusive outcomes in $\mathcal{Y}$ must sum to exactly $1$. Because the model's outputs are strictly normalized (as ensured by the softmax function in Equation \ref{eq:culprit_softmax}), we have:
\begin{equation}
    \sum_{y \in \mathcal{Y}} p^{(1)}_y = 1, \quad \sum_{y \in \mathcal{Y}} p^{(2)}_y = 1, \quad \text{and} \quad \sum_{y \in \mathcal{Y}} p^{(1,2)}_y = 1
\end{equation}

Assume for the sake of contradiction that the joint distribution is perfectly additive, meaning $p^{(1,2)}_y = p^{(1)}_y + p^{(2)}_y$ for all $y \in \mathcal{Y}$. Summing both sides over the entire prediction space yields:
\begin{align}
    \sum_{y \in \mathcal{Y}} p^{(1,2)}_y &= \sum_{y \in \mathcal{Y}} \left( p^{(1)}_y + p^{(2)}_y \right) \\
    \sum_{y \in \mathcal{Y}} p^{(1,2)}_y &= \sum_{y \in \mathcal{Y}} p^{(1)}_y + \sum_{y \in \mathcal{Y}} p^{(2)}_y \\
    1 &= 1 + 1 \\
    1 &\neq 2
\end{align}

This contradiction demonstrates that summing the resulting probability distributions of individual interventions inherently violates the fundamental bounds of probability. Therefore, $p^{(1,2)} \neq p^{(1)} + p^{(2)}$, meaning the joint intervention is strictly non-additive in probability space. \hfill $\blacksquare$

Since the joint effect is the basis for calculating the effects means that its for certain that concepts that are independ are "punished" for being in a coalition. Or mathmatically, $Make formula here$. This is no good basis for further work, so an adjusted is needed.

However removing the softmax function and changing the bayesian sampling process would disrupt backwards compatibility (\ref{des:backwards_compatibility}). Therefore the ideal way is too not remove the culprit but to work around the unwanted behavior.

Therefore the following plan: using the effect estimation with the original bayesian regression by Matton et al and then normalizing the value for intervention groups with multiple items to counter this described phenomenon. More specifically, a intervention-specific \emph{Payout Coefficient} is calculated and then ...

\subsubsection{Calculating Payout Coeefficient}
As the first step of this result normaization, the Payout coefficient is being determined for each intervention. Its the ratio of the joint effect all concept modifications have together, compared to the sum of their indivitual atomic contributions. Formally its defined with $T$ being the set of concepts to modify in this treatment.
\begin{equation}
\label{eq:payout_coefficient}
\gamma_p(x, T) = \frac{
D_{\mathrm{KL}}
\left(
\mathbb{P}_M(Y \mid x_{t \rightarrow t'})
\;\middle\|\;
\mathbb{P}_M(Y \mid x)
\right)
}{\sum{t in T}
D_{\mathrm{KL}}
\left(
\mathbb{P}_M(Y \mid x_{t \rightarrow t'})
\;\middle\|\;
\mathbb{P}_M(Y \mid x)
\right)}
\end{equation}

\subsubsection{Calculating adjusted KL Div}
Having a value $\gamma_p(x, T)$ enables to know the degree to which the joint intervention of the concepts differ from the individual impact they'd have on their own. This value then is used to adjust the KL divergence an intervention brings.

However, not the $\gamma_p(x, T)$ of the own intervention $T$ is taken, because this would make no sense. Rather, its of interest how the effect of this intervention is compared to other interventions in the dataset. Therfore a mean of $5$ Payout Coefficients is used, where datapoints are drawn by three additional criteria. However the prime criteria which cannot be skipped and is always the case is that the number of intervened concepts in the intervention always matches. So if we are looking for datapoints for $T$ and will draw from $U$ then $|T| = |U|$ otherwise we have to skip it. If one the three criteria can not supply all datapoints still needed, it goes down to the next cathegory to have a wider offering of values.

Descending priority... The mathmatical notation here is a total disaster.
\begin{enumerate}
\item Other Intervention is in the same correlaiton group. This shows that $U$ is also the result of the same correlation subgroup $s$ as determined by the oracle, sharing major similarities. Mathmatically: $|T| = |U| and s yields T and s yields U$
\item Other Intervention is of the same question $x$. $|T| = |U| and x \in T and x \in U$
\item Other Intervention is of same dataset $X$. $|T| = |U| and x \in X \in T and x \in X \in U$
\end{enumerate}

Then when the reference Payout Coefficients are gathered, their average is divided agains the origin KL div for the intervention, adjusting its effect.

\begin{equation}
\label{eq:adjusted_kl_div}
D_{\mathrm{KL_{adj}} =
D_{\mathrm{KL}
\left(
\mathbb{P}_M(Y \mid x_{t \rightarrow t'})
\;\middle\|\;
\mathbb{P}_M(Y \mid x)
\right) * \gamma_p(x, T)
}
\end{equation}

In short this normalization process ensures that the effect of all joint interventions on average is exact that of the axiomatic interventions on a dataset-level. This assumes that concepts grouped in the oracle are independent or at least synergistic and redundant effects each cancel each out on the dataset level, but first ensures entire backwards-compatibility \ref{des:backwards_compatibility} and also it ensures that concepts are not systemetcially "punished" for being put in a correlation group with others (as discuessed in the motivation part of this chapter).

SHOW THE AVERAGE EFFECT OF THIS AS DESCRIPED IN A MATHMATICAL MANNER!!!

Now the adjusted kl div formula from \ref{eq:adjusted_kl_div} is just inserted into the value function, which itself is used for \emph{Causal concept} estimation.

TODO!!! SHOWING THE UPDATED value FUNCTION HERE WITH USIGN ADJUSTED KL

\subsection{Remaining evaluationn pipeline}
Explain that now the alingment between EE and CE is still measured. However the CE definiton got changed to be based on the adjusted kldiv for each concept. Also explain why EE did not get reworked in any way.

\section{Analysis}
\subsection{Differences between extention and original}
With beatiful graphs
\subsection{Check: are the Desirada fufilled?}
With proof sketch for backwards compatibility (see presentation)
\subsection{How is the metric extention expected to influence results}
